\documentclass[10pt,a4paper]{amsart}
\usepackage{amsmath,amssymb,amsthm,graphicx}
% For commutative diagrams we use Xy-pic
\usepackage[all]{xy}

% We use 2cell for 2-commutative diagrams.
\xyoption{2cell}
% Theorem environments.
%
\theoremstyle{plain}
\newtheorem{theorem}[subsection]{Theorem}
\newtheorem{proposition}[subsection]{Proposition}
\newtheorem{lemma}[subsection]{Lemma}

\theoremstyle{definition}
\newtheorem{definition}[subsection]{Definition}
\newtheorem{example}[subsection]{Example}
\newtheorem{exercise}[subsection]{Exercise}
\newtheorem{situation}[subsection]{Situation}

\theoremstyle{remark}
\newtheorem{remark}[subsection]{Remark}
\newtheorem{remarks}[subsection]{Remarks}

\numberwithin{equation}{subsection}

% Macros
%
\def\lim{\mathop{\mathrm{lim}}\nolimits}
\def\colim{\mathop{\mathrm{colim}}\nolimits}
\def\Spec{\mathop{\mathrm{Spec}}}
\def\Hom{\mathop{\mathrm{Hom}}\nolimits}
\def\Ext{\mathop{\mathrm{Ext}}\nolimits}
\def\SheafHom{\mathop{\mathcal{H}\!\mathit{om}}\nolimits}
\def\SheafExt{\mathop{\mathcal{E}\!\mathit{xt}}\nolimits}
\def\Sch{\mathit{Sch}}
\def\Mor{\mathop{\mathrm{Mor}}\nolimits}
\def\Ob{\mathop{\mathrm{Ob}}\nolimits}
\def\Sh{\mathop{\mathit{Sh}}\nolimits}
\def\NL{\mathop{N\!L}\nolimits}
\def\CH{\mathop{\mathrm{CH}}\nolimits}
\def\proetale{{pro\text{-}\acute{e}tale}}
\def\etale{{\acute{e}tale}}
\def\QCoh{\mathit{QCoh}}
\def\Ker{\mathop{\mathrm{Ker}}}
\def\Im{\mathop{\mathrm{Im}}}
\def\Coker{\mathop{\mathrm{Coker}}}
\def\Coim{\mathop{\mathrm{Coim}}}

% Boxtimes
%
\DeclareMathSymbol{\boxtimes}{\mathbin}{AMSa}{"02}

%
% Macros for moduli stacks/spaces
%
\def\QCohstack{\mathcal{QC}\!\mathit{oh}}
\def\Cohstack{\mathcal{C}\!\mathit{oh}}
\def\Spacesstack{\mathcal{S}\!\mathit{paces}}
\def\Quotfunctor{\mathrm{Quot}}
\def\Hilbfunctor{\mathrm{Hilb}}
\def\Curvesstack{\mathcal{C}\!\mathit{urves}}
\def\Polarizedstack{\mathcal{P}\!\mathit{olarized}}
\def\Complexesstack{\mathcal{C}\!\mathit{omplexes}}
% \Pic is the operator that assigns to X its picard group, usage \Pic(X)
% \Picardstack_{X/B} denotes the Picard stack of X over B
% \Picardfunctor_{X/B} denotes the Picard functor of X over B
\def\Pic{\mathop{\mathrm{Pic}}\nolimits}
\def\Picardstack{\mathcal{P}\!\mathit{ic}}
\def\Picardfunctor{\mathrm{Pic}}
\def\Deformationcategory{\mathcal{D}\!\mathit{ef}}
\usepackage{cleveref}
\crefname{originallemma}{original lemma}{original lemmas}
\Crefname{originallemma}{Original lemma}{Original lemmas}
\makeatletter
\newlabel{sites-modules-definition-2-morphism-ringed-topoi}{{04IC}{source}{}{sourcepointer-1}{}}
\newlabel{sites-modules-definition-2-morphism-ringed-topoi@cref}{{[originallemma][0][]04IC}{[1][0][]source}{}{}{}}
\makeatother
\crefname{originallemma}{original lemma}{original lemmas}
\Crefname{originallemma}{Original lemma}{Original lemmas}
\makeatletter
\newlabel{algebra-lemma-etale-standard-smooth}{{00U9}{source}{}{sourcepointer-2}{}}
\newlabel{algebra-lemma-etale-standard-smooth@cref}{{[originallemma][0][]00U9}{[1][0][]source}{}{}{}}
\makeatother
\crefname{originaltheorem}{original theorem}{original theorems}
\Crefname{originaltheorem}{Original theorem}{Original theorems}
\makeatletter
\newlabel{theorem-exactness-stalks}{{03PU}{source}{}{sourcepointer-3}{}}
\newlabel{theorem-exactness-stalks@cref}{{[originaltheorem][0][]03PU}{[1][0][]source}{}{}{}}
\makeatother
\crefname{originallemma}{original lemma}{original lemmas}
\Crefname{originallemma}{Original lemma}{Original lemmas}
\makeatletter
\newlabel{lemma-points-small-etale-site}{{04HU}{source}{}{sourcepointer-4}{}}
\newlabel{lemma-points-small-etale-site@cref}{{[originallemma][0][]04HU}{[1][0][]source}{}{}{}}
\makeatother
\crefname{originallemma}{original lemma}{original lemmas}
\Crefname{originallemma}{Original lemma}{Original lemmas}
\makeatletter
\newlabel{algebra-lemma-strictly-henselian-solutions}{{04GX}{source}{}{sourcepointer-5}{}}
\newlabel{algebra-lemma-strictly-henselian-solutions@cref}{{[originallemma][0][]04GX}{[1][0][]source}{}{}{}}
\makeatother
\crefname{originallemma}{original lemma}{original lemmas}
\Crefname{originallemma}{Original lemma}{Original lemmas}
\makeatletter
\newlabel{topologies-lemma-alternative}{{04HR}{source}{}{sourcepointer-6}{}}
\newlabel{topologies-lemma-alternative@cref}{{[originallemma][0][]04HR}{[1][0][]source}{}{}{}}
\makeatother
\title{Morphisms of locally ringed small etale topoi of affine schemes come from morphisms of schemes}
\author{The Stacks Project Authors}
\date{stacks-project:etale-cohomology.tex:local-snapshot-20261011; GNU FDL 1.2 or later}
\begin{document}
\maketitle
\noindent\textit{Editorial note.} All delivered mathematical text is the author's own, copied from the pinned Stacks Project chapter file etale-cohomology.tex; no statement or proof step was rewritten, reordered or removed.\par\medskip
\noindent\textit{Editorial note.} The proof is the author's complete argument for existence and uniqueness: the affine computations, the stalk check, the functoriality, the identification of pushforwards and the comparison of structural maps are all inside the counted span.\par\medskip
\noindent\textit{Editorial note.} Every labelled reference occurring inside the counted spans points to a chapter of the Project that is not part of this card's source tree. Each is carried in assembly.external\_references as a non-load-bearing navigation pointer whose source location is the invocation inside the counted span, so no load-bearing mathematics is replaced by a stub.\par\medskip
\noindent\textit{Editorial note.} The checked chapter list for same-name collisions is the set of stacks chapters staged in this repository: chow, cohomology, divisors, etale-cohomology, flat, more-algebra, more-morphisms, morphisms, spaces-more-morphisms, varieties together with this card's chapter etale-cohomology.\par\medskip
\noindent\textit{Editorial note.} The preamble is the chapter's own native preamble excerpt already used for this source: the Xy-pic setup (native preamble lines 14-18) and the native theorem declarations with the mathematical macros (native preamble lines 154-217).\par\medskip
\noindent\textit{Allowed standard prerequisites (editorial).}\begin{itemize}
\item Small etale topos of a scheme with its structure sheaf and the associated locally ringed topos. Same chapter, the topos Sh(X\_etale) together with the structure sheaf, and the morphism (f\_\{small\}, f\_\{small\}\textasciicircum{}sharp) attached to a morphism of schemes. Supplies the objects compared by the asserted 2-isomorphism and the target of the construction.
\item Morphism of locally ringed topoi and 2-isomorphism of ringed topoi. Modules on Sites, the definition of a morphism of ringed topoi, of a 2-morphism between them, and of the locally ringed condition. Supplies the language in which the statement is phrased and the compatibility that the last paragraph of the proof checks.
\item Sections of the structure sheaf as maps h\_V -> O\_X and the Cech-style computation of maps out of representable sheaves. Same chapter, the description of the sheaf h\_V represented by an object of the small etale site and of morphisms from it, together with the fact that the small etale site has enough points. Supplies the computation of the isomorphism g\textasciicircum{}\{-1\}h\_V = h\_\{X x\_Y V\} and its functoriality, which produce the morphism of schemes f.
\item Uniqueness of the morphism of schemes underlying a morphism of topoi. Same chapter, the result that the small morphism of ringed topoi determines the morphism of schemes. Supplies the uniqueness clause of the statement.
\end{itemize}\medskip
\ifdefined\hypertarget\hypertarget{sourcepointer-1}{}\fi
\noindent\textit{Source reference.} Original lemma 04IC: sources/stacks-etale-cohomology/pinned-original-source.tex, source lines 5361--5361. Navigation pointer invoked at proof line 5361; the counted span contains the pointer only and the cited result is not reproduced inside it, so the pointer is registered as a non-load-bearing navigation stub in the established form for Stacks material. No bibliography citation is involved.\par\medskip
\ifdefined\hypertarget\hypertarget{sourcepointer-2}{}\fi
\noindent\textit{Source reference.} Original lemma 00U9: sources/stacks-etale-cohomology/pinned-original-source.tex, source lines 5382--5382. Navigation pointer invoked at proof line 5382; the counted span contains the pointer only and the cited result is not reproduced inside it, so the pointer is registered as a non-load-bearing navigation stub in the established form for Stacks material. No bibliography citation is involved.\par\medskip
\ifdefined\hypertarget\hypertarget{sourcepointer-3}{}\fi
\noindent\textit{Source reference.} Original theorem 03PU: sources/stacks-etale-cohomology/pinned-original-source.tex, source lines 3546--3546. Navigation pointer invoked at proof line 5428; the counted span contains the pointer only and the cited result is not reproduced inside it, so the pointer is registered as a non-load-bearing navigation stub in the established form for Stacks material. No bibliography citation is involved.\par\medskip
\ifdefined\hypertarget\hypertarget{sourcepointer-4}{}\fi
\noindent\textit{Source reference.} Original lemma 04HU: sources/stacks-etale-cohomology/pinned-original-source.tex, source lines 3610--3610. Navigation pointer invoked at proof line 5433; the counted span contains the pointer only and the cited result is not reproduced inside it, so the pointer is registered as a non-load-bearing navigation stub in the established form for Stacks material. No bibliography citation is involved.\par\medskip
\ifdefined\hypertarget\hypertarget{sourcepointer-5}{}\fi
\noindent\textit{Source reference.} Original lemma 04GX: sources/stacks-etale-cohomology/pinned-original-source.tex, source lines 5448--5448. Navigation pointer invoked at proof line 5448; the counted span contains the pointer only and the cited result is not reproduced inside it, so the pointer is registered as a non-load-bearing navigation stub in the established form for Stacks material. No bibliography citation is involved.\par\medskip
\ifdefined\hypertarget\hypertarget{sourcepointer-6}{}\fi
\noindent\textit{Source reference.} Original lemma 04HR: sources/stacks-etale-cohomology/pinned-original-source.tex, source lines 5290--5290. Navigation pointer invoked at proof line 5290; the counted span contains the pointer only and the cited result is not reproduced inside it, so the pointer is registered as a non-load-bearing navigation stub in the established form for Stacks material. No bibliography citation is involved.\par\medskip
\begin{lemma}
\label{lemma-morphism-ringed-etale-topoi-affines}
Let $X$, $Y$ be affine schemes.
Let
$$
(g, g^\#) :
(\Sh(X_\etale), \mathcal{O}_X)
\longrightarrow
(\Sh(Y_\etale), \mathcal{O}_Y)
$$
be a morphism of locally ringed topoi. Then there exists a
unique morphism of schemes $f : X \to Y$ such that
$(g, g^\#)$ is $2$-isomorphic to $(f_{small}, f_{small}^\sharp)$,
see
Modules on Sites,
Definition \ref{sites-modules-definition-2-morphism-ringed-topoi}.
\end{lemma}

\begin{proof}
In this proof we write $\mathcal{O}_X$ for the structure sheaf
of the small \'etale site $X_\etale$, and similarly for
$\mathcal{O}_Y$. Say $Y = \Spec(B)$ and $X = \Spec(A)$. Since
$B = \Gamma(Y_\etale, \mathcal{O}_Y)$,
$A = \Gamma(X_\etale, \mathcal{O}_X)$
we see that $g^\sharp$ induces a ring map $\varphi : B \to A$.
Let $f = \Spec(\varphi) : X \to Y$ be the corresponding morphism
of affine schemes. We will show this $f$ does the job.

\medskip\noindent
Let $V \to Y$ be an affine scheme \'etale over $Y$. Thus we may write
$V = \Spec(C)$ with $C$ an \'etale $B$-algebra. We can write
$$
C = B[x_1, \ldots, x_n]/(P_1, \ldots, P_n)
$$
with $P_i$ polynomials such that $\Delta = \det(\partial P_i/ \partial x_j)$
is invertible in $C$, see for example
Algebra, Lemma \ref{algebra-lemma-etale-standard-smooth}.
If $T$ is a scheme over $Y$, then a $T$-valued point of $V$ is given by
$n$ sections of $\Gamma(T, \mathcal{O}_T)$ which satisfy the polynomial
equations $P_1 = 0, \ldots, P_n = 0$. In other words, the sheaf $h_V$
on $Y_\etale$ is the equalizer of the two maps
$$
\xymatrix{
\prod\nolimits_{i = 1, \ldots, n} \mathcal{O}_Y
\ar@<1ex>[r]^a \ar@<-1ex>[r]_b &
\prod\nolimits_{j = 1, \ldots, n} \mathcal{O}_Y
}
$$
where $b(h_1, \ldots, h_n) = 0$ and
$a(h_1, \ldots, h_n) =
(P_1(h_1, \ldots, h_n), \ldots, P_n(h_1, \ldots, h_n))$.
Since $g^{-1}$ is exact we conclude that the top row of the
following solid commutative diagram is an equalizer diagram as well:
$$
\xymatrix{
g^{-1}h_V \ar[r] \ar@{..>}[d] &
\prod\nolimits_{i = 1, \ldots, n} g^{-1}\mathcal{O}_Y
\ar@<1ex>[r]^{g^{-1}a} \ar@<-1ex>[r]_{g^{-1}b} \ar[d]^{\prod g^\sharp} &
\prod\nolimits_{j = 1, \ldots, n} g^{-1}\mathcal{O}_Y \ar[d]^{\prod g^\sharp}\\
h_{X \times_Y V} \ar[r] &
\prod\nolimits_{i = 1, \ldots, n} \mathcal{O}_X
\ar@<1ex>[r]^{a'} \ar@<-1ex>[r]_{b'} &
\prod\nolimits_{j = 1, \ldots, n} \mathcal{O}_X  \\
}
$$
Here $b'$ is the zero map and $a'$ is the map defined by the
images $P'_i = \varphi(P_i) \in A[x_1, \ldots, x_n]$ via the same
rule
$a'(h_1, \ldots, h_n) =
(P'_1(h_1, \ldots, h_n), \ldots, P'_n(h_1, \ldots, h_n))$.
that $a$ was defined by. The commutativity of the diagram follows from
the fact that $\varphi = g^\sharp$ on global sections. The lower
row is an equalizer diagram also, by exactly the same arguments as
before since $X \times_Y V$ is the affine scheme
$\Spec(A \otimes_B C)$ and
$A \otimes_B C = A[x_1, \ldots, x_n]/(P'_1, \ldots, P'_n)$.
Thus we obtain a unique dotted arrow
$g^{-1}h_V \to h_{X \times_Y V}$ fitting into the diagram

\medskip\noindent
We claim that the map of sheaves $g^{-1}h_V \to h_{X \times_Y V}$
is an isomorphism. Since the small \'etale site of $X$ has enough points
(Theorem \ref{theorem-exactness-stalks})
it suffices to prove this on stalks. Hence let $\overline{x}$ be a
geometric point of $X$, and denote $p$ the associate point of the
small \'etale topos of $X$. Set $q = g \circ p$. This is a point of
the small \'etale topos of $Y$. By
Lemma \ref{lemma-points-small-etale-site}
we see that $q$ corresponds to a geometric point $\overline{y}$ of
$Y$. Consider the map of stalks
$$
(g^\sharp)_p :
(\mathcal{O}_Y)_{\overline{y}} =
\mathcal{O}_{Y, q} =
(g^{-1}\mathcal{O}_Y)_p
\longrightarrow
\mathcal{O}_{X, p} =
(\mathcal{O}_X)_{\overline{x}}
$$
Since $(g, g^\sharp)$ is a morphism of {\it locally} ringed topoi
$(g^\sharp)_p$ is a local ring homomorphism of strictly henselian
local rings. Applying localization to the big commutative diagram above and
Algebra, Lemma \ref{algebra-lemma-strictly-henselian-solutions}
we conclude that $(g^{-1}h_V)_p \to (h_{X \times_Y V})_p$ is an isomorphism
as desired.

\medskip\noindent
We claim that the isomorphisms $g^{-1}h_V \to h_{X \times_Y V}$ are
functorial. Namely, suppose that $V_1 \to V_2$ is a morphism of affine
schemes \'etale over $Y$. Write
$V_i = \Spec(C_i)$ with
$$
C_i = B[x_{i, 1}, \ldots, x_{i, n_i}]/(P_{i, 1}, \ldots, P_{i, n_i})
$$
The morphism $V_1 \to V_2$ is given by a $B$-algebra map $C_2 \to C_1$
which in turn is given by some polynomials
$Q_j \in B[x_{1, 1}, \ldots, x_{1, n_1}]$ for $j = 1, \ldots, n_2$.
Then it is an easy matter to show that the diagram of sheaves
$$
\xymatrix{
h_{V_1} \ar[d] \ar[r] & \prod_{i = 1, \ldots, n_1} \mathcal{O}_Y
\ar[d]^{Q_1, \ldots, Q_{n_2}}\\
h_{V_2} \ar[r] & \prod_{i = 1, \ldots, n_2} \mathcal{O}_Y
}
$$
is commutative, and pulling back to $X_\etale$ we obtain the
solid commutative diagram
$$
\xymatrix{
g^{-1}h_{V_1} \ar@{..>}[dd] \ar[rrd] \ar[r] &
\prod_{i = 1, \ldots, n_1} g^{-1}\mathcal{O}_Y
\ar[dd]^{g^\sharp}
\ar[rrd]^{Q_1, \ldots, Q_{n_2}} \\
& & g^{-1}h_{V_2} \ar@{..>}[dd] \ar[r] &
\prod_{i = 1, \ldots, n_2} g^{-1}\mathcal{O}_Y
\ar[dd]^{g^\sharp} \\
h_{X \times_Y V_1} \ar[r] \ar[rrd] &
\prod\nolimits_{i = 1, \ldots, n_1} \mathcal{O}_X
\ar[rrd]^{Q'_1, \ldots, Q'_{n_2}} \\
& & h_{X \times_Y V_2} \ar[r] &
\prod\nolimits_{i = 1, \ldots, n_2} \mathcal{O}_X
}
$$
where $Q'_j \in A[x_{1, 1}, \ldots, x_{1, n_1}]$ is the image of
$Q_j$ via $\varphi$. Since the dotted arrows exist, make the
two squares commute, and the horizontal arrows are injective
we see that the whole diagram commutes. This proves functoriality
(and also that the construction of $g^{-1}h_V \to h_{X \times_Y V}$
is independent of the choice of the presentation, although we
strictly speaking do not need to show this).

\medskip\noindent
At this point we are able to show that $f_{small, *} \cong g_*$.
Namely, let $\mathcal{F}$ be a sheaf on $X_\etale$. For every
$V \in \Ob(X_\etale)$ affine we have
\begin{align*}
(g_*\mathcal{F})(V)
& =
\Mor_{\Sh(Y_\etale)}(h_V, g_*\mathcal{F}) \\
& =
\Mor_{\Sh(X_\etale)}(g^{-1}h_V, \mathcal{F}) \\
& =
\Mor_{\Sh(X_\etale)}(h_{X \times_Y V}, \mathcal{F}) \\
& =
\mathcal{F}(X \times_Y V) \\
& =
f_{small, *}\mathcal{F}(V)
\end{align*}
where in the third equality we use the isomorphism
$g^{-1}h_V \cong h_{X \times_Y V}$ constructed above. These isomorphisms
are clearly functorial in $\mathcal{F}$ and functorial in $V$
as the isomorphisms $g^{-1}h_V \cong h_{X \times_Y V}$ are functorial.
Now any sheaf on $Y_\etale$ is determined by the restriction
to the subcategory of affine schemes
(Topologies, Lemma \ref{topologies-lemma-alternative}),
and hence we obtain an isomorphism of functors $f_{small, *} \cong g_*$
as desired.

\medskip\noindent
Finally, we have to check that, via the isomorphism
$f_{small, *} \cong g_*$ above, the maps $f_{small}^\sharp$ and
$g^\sharp$ agree. By construction this is already the case for the
global sections of $\mathcal{O}_Y$, i.e., for the elements of $B$.
We only need to check the result on
sections over an affine $V$ \'etale over $Y$ (by
Topologies, Lemma \ref{topologies-lemma-alternative}
again). Writing
$V = \Spec(C)$, $C = B[x_i]/(P_j)$ as before it suffices
to check that the coordinate functions $x_i$ are mapped to
the same sections of $\mathcal{O}_X$ over $X \times_Y V$.
And this is exactly what it means that the diagram
$$
\xymatrix{
g^{-1}h_V \ar[r] \ar@{..>}[d] &
\prod\nolimits_{i = 1, \ldots, n} g^{-1}\mathcal{O}_Y
\ar[d]^{\prod g^\sharp} \\
h_{X \times_Y V} \ar[r] &
\prod\nolimits_{i = 1, \ldots, n} \mathcal{O}_X
}
$$
commutes. Thus the lemma is proved.
\end{proof}

\end{document}
